\section{Adding the Trakr Flat Locomotion Task to \texttt{unitree\_rl\_lab}}

The Trakr flat locomotion task was added as a new Isaac Lab extension package inside the existing
\texttt{unitree\_rl\_lab} style repository. Instead of modifying the Unitree Go2 task in-place, I created
a separate package named \texttt{trakr\_rl} under \texttt{source/trakr\_rl}. This kept the Trakr robot
asset, task registration, locomotion configuration, reward terms, and helper scripts isolated from the
original Unitree environments while still reusing the same manager-based reinforcement learning pipeline,
RSL-RL training scripts, and deployment/export flow.

\subsection{Extension and Task Registration}

The Trakr task is exposed to Isaac Lab through Gymnasium registration. The package initializer
\texttt{trakr\_rl/tasks/\_\_init\_\_.py} imports all task subpackages using Isaac Lab's package import
utility, which causes the robot task registration files to run when the extension is loaded. The flat
locomotion task is registered in
\texttt{trakr\_rl/tasks/locomotion/robots/trakr/\_\_init\_\_.py} with the environment id
\texttt{Trakr-Velocity}. This registration maps the id to Isaac Lab's
\texttt{ManagerBasedRLEnv}, points the training configuration to \texttt{RobotEnvCfg}, points the play
configuration to \texttt{RobotPlayEnvCfg}, and uses the shared PPO runner configuration in
\texttt{trakr\_rl.tasks.locomotion.agents.rsl\_rl\_ppo\_cfg}.

This design makes the task discoverable by the existing command-line workflow. After installing the
extension with \texttt{./trakr\_rl.sh --install}, the environment can be trained and played using the same
script structure as the Unitree tasks:

\begin{verbatim}
./trakr_rl.sh --train --task Trakr-Velocity
./trakr_rl.sh --play  --task Trakr-Velocity
\end{verbatim}

\subsection{Trakr Robot Asset Configuration}

The robot description is defined in \texttt{trakr\_rl/assets/robots/trakr.py}. I added a
\texttt{TrakrArticulationCfg} and used a USD-based spawn configuration for the Trakr model. The
configuration points to the inspected \texttt{trakr\_imu.usd} asset, enables contact sensors, and sets the
rigid body and articulation properties required by Isaac Lab. The rigid body parameters were adapted from
the earlier Trakr legged RL setup, while the articulation solver settings follow the Unitree Go2-style
configuration.

The Trakr-specific initial state and actuator model were also defined in this file. The default base
height is set to \texttt{0.255 m}, all joints start from a neutral standing pose, and a DC motor actuator
configuration is applied to all adduction, hip, and knee joints. The motor limits are configured with
\texttt{23.5 Nm} effort and saturation limits, \texttt{30 rad/s} velocity limit, \texttt{20.0} stiffness,
and \texttt{1.0} damping. I also listed the SDK joint ordering explicitly:
\texttt{LB}, \texttt{LF}, \texttt{RB}, and \texttt{RF}, each with adduction, hip, and knee joints. This
was necessary because the policy, deployment code, and USD articulation must agree on joint naming and
ordering.

\subsection{Flat Environment Scene}

The flat task environment is implemented in
\texttt{trakr\_rl/tasks/locomotion/robots/trakr/velocity\_env\_cfg.py}. The scene configuration replaces
the Unitree robot with \texttt{TRAKR\_CFG} and places it at the prim path
\texttt{\{ENV\_REGEX\_NS\}/trakr}. I updated all Trakr-specific prim and body names after inspecting the
USD stage in Isaac Sim. For example, the base body is referenced as \texttt{base\_link}, the height scanner
uses \texttt{\{ENV\_REGEX\_NS\}/trakr/trakr/base\_link}, and contact sensors are attached under
\texttt{\{ENV\_REGEX\_NS\}/trakr/trakr/.*}. These changes were important because the original Go2 task
uses different body names and prim paths.

Although the terrain is created through Isaac Lab's terrain generator interface, the flat locomotion task
uses only a plane sub-terrain. This keeps the training objective focused on learning a stable flat-ground
velocity tracking gait before moving to rough terrain. The scene uses a generated terrain with an
\texttt{8 m \(\times\) 8 m} tile size, a large border, friction set to \texttt{1.0}, and a marble visual
material from the Isaac Lab Nucleus assets. The environment runs \texttt{512} parallel environments for
training and \texttt{32} environments for play.

\subsection{MDP Design}

The action space is a joint position action over all Trakr joints. The action scale is \texttt{0.25}, and
the default joint offsets are used so that policy outputs are interpreted relative to the nominal standing
configuration. The policy observation group contains angular velocity, projected gravity, commanded base
velocity, relative joint position, relative joint velocity, and the previous action. A history length of
\texttt{5} is used and observation corruption is enabled, which gives the policy short temporal context
while improving robustness to sensor noise. The critic receives privileged observations, including base
linear velocity and joint effort, in addition to the policy terms.

Velocity commands are generated with a custom \texttt{UniformLevelVelocityCommandCfg}. This extends Isaac
Lab's uniform velocity command configuration with \texttt{limit\_ranges}, allowing the curriculum to start
with a small command range and gradually expand it. For the flat task, the initial command range is
\texttt{[-0.1, 0.1] m/s} in both linear directions and \texttt{[-1, 1] rad/s} in yaw. The curriculum can
increase the limits up to \texttt{[-1.0, 1.0] m/s} in the forward direction,
\texttt{[-0.4, 0.4] m/s} laterally, and \texttt{[-1.0, 1.0] rad/s} in yaw.

\subsection{Events, Rewards, and Terminations}

The event configuration randomizes physical properties at startup and reset to reduce overfitting to a
single simulation instance. At startup, rigid body friction is sampled between \texttt{0.3} and
\texttt{1.2}, restitution is sampled between \texttt{0.0} and \texttt{0.15}, and base mass is perturbed by
an additive range of \texttt{[-1.0, 3.0] kg}. During resets, the base pose is randomized in position and
yaw, joint velocities are randomized, and interval pushes are applied by setting base velocity in the
\(x\)- and \(y\)-directions.

The reward function follows the standard velocity-tracking structure from the Unitree locomotion task, but
the weights and additional penalties were adjusted for Trakr. Positive reward terms track commanded
linear and yaw velocity. Penalties discourage vertical base velocity, roll and pitch instability, excessive
joint velocity, acceleration, torque, action rate, energy use, joint limit violations, foot sliding,
uneven air time, and undesired contacts on non-foot bodies. I added custom reward utilities in
\texttt{trakr\_rl/tasks/locomotion/mdp/rewards.py}, including energy usage, roll penalty, pitch-rate
penalty, air-time variance, and foot-height terms. These terms helped shape the learned gait toward stable
flat locomotion while avoiding unstable body oscillations.

Episodes terminate on timeout, base contact, or excessive body orientation. The base-contact termination
uses \texttt{base\_link}, matching the Trakr USD body name, and the bad-orientation termination uses a
limit angle of \texttt{0.8 rad}. This prevents failed rollovers or body collisions from continuing to
contribute poor trajectories to the PPO update.

\subsection{Training Configuration}

The PPO runner is defined in
\texttt{trakr\_rl/tasks/locomotion/agents/rsl\_rl\_ppo\_cfg.py}. I reused the RSL-RL on-policy runner
interface with \texttt{32} steps per environment, a save interval of \texttt{100}, and actor--critic
networks of size \texttt{[256, 128, 64]} using ELU activations. The PPO algorithm uses a clipped objective
with \texttt{0.2} clipping, \texttt{0.01} entropy coefficient, \texttt{0.99} discount factor,
\texttt{0.95} GAE parameter, adaptive learning rate, and \texttt{1e-3} initial learning rate. Trained
checkpoints and exported policies are stored under \texttt{logs/rsl\_rl/trakr\_velocity}, confirming that
the flat Trakr task was integrated into the same training and export workflow as the original
\texttt{unitree\_rl\_lab} tasks.

\begin{table}[h]
    \centering
    \caption{Main files added or modified for the Trakr flat locomotion task.}
    \begin{tabular}{p{0.38\linewidth} p{0.52\linewidth}}
        \hline
        \textbf{File} & \textbf{Purpose} \\
        \hline
        \texttt{source/trakr\_rl} & New Isaac Lab extension package for Trakr tasks. \\
        \texttt{assets/robots/trakr.py} & Defines the Trakr USD articulation, initial pose, actuators, and joint ordering. \\
        \texttt{tasks/locomotion/robots/trakr/\_\_init\_\_.py} & Registers the flat task as \texttt{Trakr-Velocity}. \\
        \texttt{tasks/locomotion/robots/trakr/velocity\_env\_cfg.py} & Defines the flat scene, observations, actions, commands, rewards, events, curriculum, and terminations. \\
        \texttt{tasks/locomotion/mdp/rewards.py} & Adds Trakr-specific reward and penalty terms. \\
        \texttt{tasks/locomotion/mdp/curriculums.py} & Adds command range curriculum logic. \\
        \texttt{tasks/locomotion/mdp/commands/velocity\_command.py} & Extends velocity command configuration with curriculum limit ranges. \\
        \texttt{tasks/locomotion/agents/rsl\_rl\_ppo\_cfg.py} & Defines the PPO runner and actor--critic settings. \\
        \texttt{trakr\_rl.sh} & Provides install, list, train, and play entry points for the Trakr extension. \\
        \hline
    \end{tabular}
    \label{tab:trakr_flat_task_files}
\end{table}

Overall, the Trakr flat locomotion task was integrated by following the same modular structure as the
Unitree locomotion environments: a registered Gymnasium task id, a robot articulation asset, a
manager-based environment configuration, custom MDP terms, and an RSL-RL PPO runner. The main adaptation
work was replacing Go2-specific robot assumptions with Trakr USD prim paths, body names, joint names,
actuator parameters, and reward shaping suitable for the Trakr morphology.
